\subsection{Examples of approximation spaces}

Recall that every Hilbert space is an approximation space (III, p. 15, prop. 10). This section gives other examples.

If X is a locally compact space, we denote by \(\mathcal{C}_{0}(\mathrm{X};\mathrm{K})\) , or simply \(\mathcal{C}_{0}(\mathrm{X})\) , the Banach space of continuous functions from X to K vanishing at infinity, equipped with the norm defined by \(\|f\|=\sup_{x\in\mathrm{X}}|f(x)|\) for \(f\in\mathcal{C}_{0}(\mathrm{X})\) . When X is compact, this space coincides with the space \(\mathcal{C}(\mathrm{X})\) of continuous functions from X to K.

Lemma 5. --- Let X be a compact topological space, F a finite subset of \(\mathcal{C}(\mathrm{X})\) and \(\varepsilon > 0\) a real number.

There exist a finite subset \(X_{0} \subset X\) and a linear map \(u: \mathcal{C}(X_{0}) \to \mathcal{C}(X)\) of norm \(\leqslant 1\) such that \(\|u(f|X_{0}) - f\| \leqslant \varepsilon\) for all \(f \in F\) .

Since the set F is an equicontinuous subset of \(\mathcal{C}(\mathrm{X})\) there exists a finite covering \((\mathrm{U}_{i})_{i\in\mathrm{I}}\) of X such that, for every \(i\in I\) and every \(f\in F\) the diameter of \(f(\mathrm{U}_{i})\) is \(\leqslant\varepsilon\) . For each \(i\in I\), choose a point \(x_{i}\) of \(U_{i}\) , and denote by \(X_{0}\) the finite set of \(x_{i}\) for \(i\in I\) .

Let \((\varphi_{i})_{i\in\mathrm{I}}\) be a continuous partition of unity subordinate to \((\mathrm{U}_{i})_{i\in\mathrm{I}}\) . Define a linear map u from \(\mathcal{C}(\mathrm{X}_{0})\) into \(\mathcal{C}(\mathrm{X})\) by setting

\[
u (g) = \sum_ {i \in \mathrm{I}} g (x _ {i}) \varphi_ {i}
\]

for all \(g \in \mathcal{C}(X_{0})\) . The linear map u has norm \(\leqslant 1\) and satisfies \(\|u(f|X_{0}) - f\| \leqslant \varepsilon\) for all f in F, hence the assertion.

PROPOSITION 13. --- Let X be a locally compact topological space. The Banach space \(\mathcal{C}_{0}(\mathrm{X})\) is an approximation space.

Suppose first that X is compact. Let \(A \subset \mathcal{L}^{\mathrm{f}}(\mathcal{C}(\mathrm{X}))\) be the set of maps of the form \(f \mapsto u(f|X_{0})\) , where \(X_{0}\) is a finite subset of X and \(u : \mathcal{C}(X_{0}) \to \mathcal{C}(X)\) is a linear map of norm \(\leqslant 1\) . The set A is equicontinuous. By Lemma 5, the identity map of \(\mathcal{C}(X)\) is adherent to A for the topology of pointwise convergence on \(\mathcal{C}(X)\) . The proposition then follows from Remark 3 of III, p. 14.

Let us pass to the general case. Let Y be the Alexandroff compactification of X and \(\omega\) the point at infinity of Y (TG, I, p. 67). Identify \(\mathcal{C}_{0}(X)\) with the set of elements of \(\mathcal{C}(Y)\) vanishing at \(\omega\) . Then \(\mathcal{C}(Y)\) is the topological direct sum of \(\mathcal{C}_{0}(X)\) and of the vector space of constant functions on Y. Since \(\mathcal{C}(Y)\) is an approximation space according to the preceding, the space \(\mathcal{C}_{0}(X)\) is an approximation space (Remark 1 of III, p. 14).

COROLLARY. --- Every commutative star algebra is an approximation space.

Indeed, such an algebra is isomorphic to the algebra of continuous functions vanishing at infinity on a locally compact space (I, p. 108, th. 1).

Remark. --- If E is an infinite-dimensional Hilbert space, the star algebra \(\mathcal{L}(\mathrm{E})\) does not have the approximation property (A. SZANKOWSKI, \(\mathcal{B}(\mathrm{H})\) does not have the approximation property, Acta Math. 147 (1981), p. 89--108).

Lemma 6. --- Let X be a locally compact topological space. Let \(\mu\) be a positive measure on X and \(p \in [1, +\infty[\) . Let F be a finite subset of \(\mathrm{L}_{\mathrm{K}}^{p}(\mathrm{X}, \mu)\) and \(\varepsilon > 0\) a real number.

There exists a finite-rank projector u of \(\mathrm{L}_{\mathrm{K}}^{p}(\mathrm{X},\mu)\) of norm \(\leqslant 1\) such that \(\|u(f)-f\|\leqslant\varepsilon\) for every f in F.

Let \(\mathcal{P}\) be the set of finite partitions \(\pi = (\mathrm{K}_1, \ldots, \mathrm{K}_n, \mathrm{H})\) of X, where \(n \geqslant 1\) is an integer and \(\mathrm{K}_1, \ldots, \mathrm{K}_n\) are integrable subsets of X with nonzero measure. For every partition \(\pi = (\mathrm{K}_1, \ldots, \mathrm{K}_n, \mathrm{H}) \in \mathcal{P}\) , one defines an endomorphism \(v_\pi\) of \(\mathcal{L}^p(\mathrm{X}, \mu)\) by setting

\[
v _ {\pi} (f) = \sum_ {i = 1} ^ {n} \frac {1}{\mu (\mathrm{K} _ {i})} \left(\int_ {\mathrm{K} _ {i}} f d \mu\right) \varphi_ {\mathrm{K} _ {i}},
\]

for \(f \in \mathcal{L}_{\mathrm{K}}^{p}(\mathrm{X}, \mu)\) , where \(\varphi_{\mathrm{K}_i}\) is the characteristic function of \(\mathrm{K}_i\) . The map \(v_\pi\) induces, by passing to quotients, a projector \(u_\pi\) of \(\mathrm{L}_{\mathrm{K}}^{p}(\mathrm{X}, \mu)\) . One verifies easily that the image of \(u_\pi\) is the space of classes of functions \(f \in \mathcal{L}_{\mathrm{K}}^{p}(\mathrm{X}, \mu)\) such that \(f\) vanishes on H and is constant on \(\mathrm{K}_i\) for \(1 \leqslant i \leqslant n\) .

Let us prove that \(\| u_{\pi}\| \leqslant 1\) . For \(f\in \mathcal{L}_{\mathrm{K}}^{p}(\mathrm{X},\mu)\) , we have

\[
\| u _ {\pi} (f) \| _ {p} ^ {p} = \sum_ {i = 1} ^ {n} \mu (\mathrm{K} _ {i}) ^ {1 - p} \left| \int_ {\mathrm{K} _ {i}} f d \mu \right| ^ {p}.
\]

By Hölder's inequality (INT, IV, p. 208, § 6, n° 4, th. 2), we have for every i the inequality

\[
\left| \int_ {\mathrm{K} _ {i}} f d \mu \right| ^ {p} \leqslant \mu (\mathrm{K} _ {i}) ^ {p - 1} \int_ {\mathrm{K} _ {i}} | f | ^ {p} d \mu ,
\]

whence \(\| u_{\pi}(f)\| _p\leqslant \| f\| _p.\)

Let \(\mathcal{E}\) be the set of classes in \(\mathrm{L}_{\mathrm{K}}^{p}(\mathrm{X},\mu)\) of integrable functions on X which take only finitely many values. Since \(\mathcal{E}\) is dense in \(\mathrm{L}_{\mathrm{K}}^{p}(\mathrm{X},\mu)\) (INT, IV, p. 162, §4, n° 10, cor. 1), there exists a finite set \(\mathrm{F}'\) of \(\mathcal{E}\) such that every element of F is at distance at most \(\varepsilon\) from an element of \(\mathrm{F}'\) . Considering the finite partition \(\pi\) formed by the sets of non-zero measure on which the map \(x\mapsto (f(x))_{f\in \mathrm{F}'}\) takes a given value, one sees that there exists an element \(\pi\) of \(\mathcal{P}\) such that \(u_{\pi}(f) = f\) for all \(f\) in \(\mathrm{F}'\) . The projector \(u_{\pi}\) therefore has the required properties.

PROPOSITION 14. --- Let X be a locally compact topological space, \(\mu\) a positive measure on X, and \(p \in [1, +\infty]\) . The space \(\mathrm{L}_{\mathrm{K}}^{p}(\mathrm{X}, \mu)\) is an approximation space.

If \(p = +\infty\) , then the space \(\mathrm{L}_{\mathbf{C}}^{\infty}(\mathrm{X},\mu)\) is a commutative star algebra (example 4 of I, p. 103), hence an approximation space (cor. of prop. 13), and the same is true of \(\mathrm{L}_{\mathbf{R}}^{\infty}(\mathrm{X},\mu)\) by remark 5 of III, p. 15.

Suppose that p is finite. Let \(\mathrm{A}\subset\mathcal{L}^{\mathrm{f}}(\mathrm{L}_{\mathrm{K}}^{p}(\mathrm{X},\mu))\) be the set of finite-rank projectors of norm \(\leqslant1\) . By lemma 6, the identity map of \(\mathrm{L}_{\mathrm{K}}^{p}(\mathrm{X},\mu)\) is adherent to A for the topology of simple convergence, and the proposition then follows from remark 3 of III, p. 14.

